% golden A — 人手序列化（ir2tikz 的对照真值）
% source: tests/golden/A/ir.json  (voltage divider with bypass cap and port)
% nets:
%   N_IN : V1.2, R1.1
%   N_MID: R1.2, R2.1, C1.1, port a(5.5,2)   <- junction (2,2)
%   GND  : V1.1, R2.2, C1.2, GND1.p
\documentclass[margin=5pt]{standalone}
\usepackage[american]{circuitikz}
\ctikzset{bipoles/length=1.0cm}
\begin{document}
\begin{circuitikz}

%% ==== [region] source ====
\draw (0,0) to[V, invert, l=$U_s$, a=$6\mathrm{V}$, name=V1] (0,4);

%% ==== [region] divider ====
\coordinate (N_MID) at (2,2);
\coordinate (GND) at (2,0);
\draw (2,4) to[R, l=$R_1$, a=$1\mathrm{k}\Omega$, name=R1] (2,2);
\draw (2,2) to[R, l=$R_2$, a=$2\mathrm{k}\Omega$, name=R2] (2,0);

%% ==== [region] output ====
\draw (4,2) to[C, l=$C_1$, a=$100\mathrm{nF}$, name=C1] (4,0);
\node[ground] at (4,0) {};

%% ==== wires ====
% net N_IN
\draw (0,4) -- (2,4);              % W1: V1.2 -> R1.1
% net GND
\draw (0,0) -- (GND) -- (4,0);     % W2: V1.1 -> R2.2 -> C1.2
% net N_MID
\draw (N_MID) -- (4,2) -- (5.5,2); % W3: R1.2 -> C1.1 -> port a

%% ==== junctions / terminals / texts ====
\node[circ] at (N_MID) {};
\node[ocirc] at (5.5,2) {};
\node[anchor=west] at (5.8,2) {$a$};
\node[anchor=south] at (5.5,2.8) {$u_o$};

\end{circuitikz}
\end{document}
